\chapter{QUADRATIC FORMS; ELEMENTARY GEOMETRY.}

The theory of quadratic forms, in its modern aspect, hardly goes back before the second half of the XVIIIth century, and, as we will see, it was developed mainly to respond to the needs of Arithmetic, Analysis and Mechanics. But the fundamental notions of this theory have in reality appeared from the beginnings of ``Euclidean'' geometry, of which they formed the frame. For this reason, its history can not be retraced without speaking, at least in a summary way, of the development of ``elementary geometry'' since antiquity. Of course, we will only be able to link up with the evolution of a few general ideas, and the reader must not expect to find here precise information on the history of such or such a particular theorem, about which it will suffice to refer back to specialised historical or didactic works. \(^{1}\) It goes without saying also, when we speak below of various possible interpretations of the same theorem in various algebraic or geometric languages, that we do not intend at all to say that these ``translations'' have at all times been as familiar as they are today; quite on the contrary, it is the principal gaol of this Note to show how, very gradually, mathematicians have become conscious of these relationships between questions often of very different aspect; we will also have to show how, doing this, they have been led to put some coherence in the mass of theorems from geometry bequeathed by the ancients, and finally to try to delimit exactly what should be understood by ``geometry''.

If one puts aside the discovery, by the Babylonians, of the formula for the solution of equations of the second degree ({[}232{]}, pp.~183-189) it is therefore under their geometric disguise that the birth of the principal concepts of the theory of quadratic forms must be noted. These appear first as squares of distances (in the plane or space of three dimensions) and the corresponding notion of ``orthogonality'' is introduced by means of the right angle, defined by Euclid as half the straight angle (Elements, Book I, Def. 10); the notions of distance and right angle being linked by the theorem of Pythagoras, the keystone of the Euclidean edifice. \(^{2}\) The idea of angle seems to have been introduced very early in Greek mathematics (which no doubt had received it from the Babylonians, broken in to the use of angles by their long astronomical experience). It is known that in the classical period, only angles less than 2 right angles are defined (the ``definition'' of Euclid is besides as vague and unusable as the ones he gives of straight line and plane); the notion of direction is not brought out, although Euclid uses (without axiom or definition) the fact that a straight line divides the plane into two regions, which he distinguishes carefully when it is necessary. \(^{3}\) At this stage, the idea of the group of plane rotations only comes to light in a very imperfect way, by addition (introduced, it also, without explanation by Euclid) of non directed angles of half-lines, which is only defined, in principle, when the sum is at most equal to two right angles. \(^{4}\) As for trigonometry, it is disdained by geometers and abandoned to surveyors and astronomers; it is these latter (Aristarchus, Hipparchus, Ptolemy especially {[}255{]}) who establish the fundamental relations between the sides and angles of a right angled triangle (plane or spherical) and draw up the first tables (they consist of tables giving the chord of the arc cut out by an angle \(\theta < \pi\) on a circle of radius r, in other words the number \(2r \sin \frac{\theta}{2}\) ; the introduction of the sine, more easily handled, is due to Hindu mathematicians of the Middle Ages); in the calculation of these tables, the formula for the addition of arcs, unknown at this time, is replaced by the equivalent use of the theorem of Ptolemy (going back perhaps to Hipparchus)

on quadrilaterals inscribed in a circle. It must be noted also that Euclid and Hero give propositions equivalent to the formula

\[
a ^ {2} = b ^ {2} + c ^ {2} - 2 b c \cos A
\]

between the sides and angles of an arbitrary plane triangle; but one can hardly see there the first appearance of the notion of a bilinear form associated with a metric form, for lack of the idea of a vector calculus which will only emerge in the XIXth century.

Displacements (or movements, the distinction between the two notions not being clear in antiquity --- nor even much later) are known to Euclid; but, for reasons of which we are ignorant, he seems to feel a clear repugnance to use it (for example in the ``case of equality of triangles'', where the impression is gained that he only uses the notion of displacement because he has not been able to formulate an appropriate axiom ({[}153 e{]}, v. I, pp.~225-227 and 249)); however, it is to the notion of displacement (rotation about an axis) that he has recourse for the definition of cones of revolution and of spheres (Elements, Book XI, def. 14 and 18), as well as Archimedes for that of quadrics of revolution. But the general idea of a transformation, applied to the whole of space, is more or less a stranger to mathematical thought before the end of the XVIIIth century;⁵ and before the XVIIth century, no trace is found either of the notion of composition of movements, nor a fortiori of the composition of displacements. That does not mean to say, of course, that the Greeks were not particularly aware of the ``regularities'' and ``symmetries'' of figures, which we now attach to the notion of groups of displacements: their theory of regular polygons, and even more that of regular polyhedra --- one of the most remarkable chapters of all their mathematics --- is there to prove the opposite.⁶

Finally, the last of the essential contributions of Greek mathematics, in the area which concerns us, is the theory of conics (as far as quadrics is concerned, the Greeks only know certain quadrics of revolution, and do not carry their study of them very far, except for the sphere). It is interesting to note here that, although the Greeks never had the idea of the fundamental principle of analytic geometry (essentially because of the lack of a manageable algebra), they commonly use, for the study of particular ``figures'', ``ordinates'' relative to two (or even more than two) axes in the plane (linked closely to the figure, which is one of the fundamental points where their method differs from that of Fermat and Descartes, whose axes are fixed independently of the figure under consideration). In particular, the first examples of conics (other than the circle) which come up in connection with the problem of the duplication of the cube, are the curves given by the equations \(y^{2}=ax, y=bx^{2}, xy=c\) (Menechme, a pupil of Eudoxus, middle of the IVth century); \(^{7}\) and it is the equation of conics (normally with reference to two oblique axes made up of a diameter and of the tangent to one of its points of contact with the curve) that is most often used in the study of problems relative to these curves (whereas the ``focal'' properties only play a very effacing role, contrary to what scholarly traditions going back only to the XIXth century would make us believe). From this vast theory, we must above all remember here the notion of conjugate diameters (already known by Archimedes), and the property that today serves as the definition of the polar of a point, given by Apollonius {[}153 b{]} when the point is exterior to the conic (the polar being therefore for him the straight line joining the points of contact of the tangents issuing from this point); from our point of view, these are two examples of ``orthogonality'' with respect to a quadratic form distinct from the metric form, but it is well understood that the link between these notions and the classical notion of perpendiculars could not possibly be conceived of at that time.

There is hardly any other progress to point out before Descartes and Fermat; but from the beginnings of analytic geometry, the algebraic theory of quadratic forms starts to emerge from its geometric crust: Fermat knows that an equation of the second degree in the plane represents a conic ({[}109{]}, v. I, pp.~100-102; French translation, v. III, pp.~84-101) and sketches analogous ideas about quadrics ({[}109{]}, v. I, pp.~111-117; French translation, v. III, pp.~102-108). With the development of analytic geometry in 2 and 3 dimensions during the course of the XVIIIth century two of the central problems of the theory appear (especially in connection with conics and quadrics): the reduction of a quadratic form to a sum of squares and the search for its ``axes'' with respect to the metric form. For conics, these two problems are too elementary to give rise to important algebraic progress; for an arbitrary number of variables, the first is solved by Lagrange in 1759, in connection with the maxima of functions of several variables ({[}191{]}, v. I, pp.~3-20). But this problem is almost immediately eclipsed by that of the search for axes, even before invariance of rank had been formulated; \(^{8}\) as for the law of inertia, it is only discovered around 1850 by Jacobi ({[}171{]}, v. III, pp.~593-598), who proves it by the same argument as nowadays, and Sylvester ({[}304{]}, v. I, pp.~378-381) who limits himself to stating it as quasi-obvious. \(^{9}\)

The problem of the reduction of a quadric to its axes already presents algebraic difficulties substantially greater than the analogous problem for conics; and Euler, who is the first to tackle it, is not in a position to prove the real character of the eigenvalues, which he admits after an attempt at justification without conviction ({[}108 a{]}, (1), v. IX, pp.~379-392). \(^{10}\) Given that this point is correctly established in about 1800 {[}140{]}, we must await Cauchy for the proof of the corresponding theorem for forms in an arbitrary number n of variables ({[}56 a{]}, (2), v. IX, pp.~174-195). It is also Cauchy who, at about the same time, proves that the characteristic equation giving the eigenvalues is invariant under all orthogonal changes of axes ({[}56 a{]}, (2), v. V, p.~252); \(^{11}\) but for n = 2 or n = 3, this invariance was intuitively ``obvious'' because of the geometric interpretation of eigenvalues by means of axes of the corresponding conic or quadric. Besides, during the course of research on this subject, the elementary symmetric functions of the eigenvalues have also occurred naturally (with various geometric interpretations, in connection notably with the theorems of Apollonius on conjugate diameters), and in particular the discriminant, which (known for a long time for n = 2 in connection with the theory of the equation of the second degree) appears for the first time for n = 3 with Euler ({[}108 a{]}, (1), v. IX, p.~382); this latter meets it in connection with the classification of quadrics (in expressing the condition for a quadric not to have a point at infinity) and does not mention its invariance with respect to orthogonal axis changes. But a bit later, with the beginnings of the arithmetical theory of quadratic forms with integral coefficients, Lagrange notes (for n = 2) a particular case of the invariance of the discriminant under linear but not orthogonal change of variables ({[}191{]}, v. III, p.~699), and Gauss establishes, for n = 3, the ``covariance'' of the discriminant for all linear transformations ({[}124 a{]}, v. I, pp.~301-302). \(^{12}\) Once the general formula for the multiplication of determinants was proved by Cauchy and Binet, the extension of the formula of Gauss to an arbitrary number of variables was immediate; it is that which, around 1845, will give the first impulsion to the general theory of invariants.

To the two notions that, amongst the Greeks, took the place of the theory of displacements --- that of movement and that of ``symmetry'' of a figure --- has just been added a third in the XVIIth and XVIIIth centuries with the problem of the changing of rectangular axes, which is substantially equivalent to this theory. Euler devotes several works to this question, attaching importance above all to obtaining parametric representations which are manageable for the formulae for changing axes. It is known what use Mechanics was going to make of the three angles which he introduces for this purpose for n = .3 ({[}108 a{]}, (1), v. IX, pp.~371-378). But he does not limit himself to that, considers in 1770 the general problem of orthogonal transformations for arbitrary n, remarks that the goal is reached here by introducing \(n(n - 1)/2\) angles as parameters, and finally, for n = 3 and n = 4, gives for rotations rational representations (as a function, respectively, of 4 homogeneous parameters and of 8 homogeneous parameters linked by a relation), which are none other than those obtained later by means of the theory of quaternions, and of which he does not indicate the origin ({[}108 a{]}, (1), v. VI, pp.~287-315). \(^{13}\)

On the other hand, Euler indicates also how to translate analytically the search for the ``symmetries'' of plane figures, and it is for this purpose that he is led to prove, substantially, that a plane displacement is a rotation, or a translation, or a translation followed by a symmetry ({[}108 a{]}, (1), v. IX, pp.~197-199). The rise of Mechanics at this time leads besides to the general study of displacements: but first of all it is only a question of ``infinitely small'' displacements tangent to continuous movements: they are apparently the only ones that occur in the researches of Torricelli, Roberval and Descartes on the composition of movements and the instantaneous centre of rotation for plane movements (cf.~p.~177). This latter is defined in a general way by Johann Bernoulli; d'Alembert in 1749, Euler the following year, extend this notion by proving the existence of an instantaneous axis of rotation for movements which leave a point fixed. The analogous theorem for finite displacements is not stated until 1775 by Euler {[}108 b{]}, in a memoir where he discovers at the same time that the determinant of a rotation is equal to 1; the following year, he proves the existence of a fixed point for plane similarities ({[}108 a{]}, (1), v. XXVI, pp.~276-285). But the work of Chasles must be awaited, starting in 1830 {[}60 a{]}, in order to obtain finally a coherent theory of finite and infinitely small displacements.

We arrive thus at what can be called the golden age of geometry, which is inserted grosso modo between the dates of the publication of the Géométrie descriptive of Monge (1795) {[}225{]} and of the ``Erlangen programme'' of F. Klein (1872) ({[}182{]}, v. I, pp.~460-497). The essential progress that we owe to this brusque renewal of geometry are the following:

\begin{enumerate}
\def\labelenumi{\Alph{enumi})}
\item
  The notion of the element at infinity (point, straight line or plane), introduced by Desargues in the XVIIth century {[}84{]}, but which is hardly expressed during the XVIIIth century except by abuse of language, is rehabilitated and systematically used by Poncelet {[}252{]} who thus makes of projective space the general framework of all geometrical phenomena.
\item
  At the same time, with Monge and especially Poncelet, the passage to complex projective geometry is effected. The notion of imaginary point, sporadically used during the XVIIIth century, is here exploited (concurrently with that of the point at infinity) to give statements independent of the ``case of the figure'' of real affine geometry. If at first the justifications brought in to support these innovations remained very awkward (especially those from the members of the school of ``synthetic'' geometry, where the use of coordinates ends up as being seen as a disgrace), what cannot be missed is recognising there, under the name of ``principle of contingent relations'' in Monge, or ``principle of continuity'' in Poncelet, the first germ of the idea of ``specialisation'' from modern algebraic geometry. \(^{14}\)
\end{enumerate}

One of the first results following on from these concepts is the remark that, in complex projective space, all the conics (respectively quadrics) which are non degenerate are of the same kind; which brings Poncelet to the discovery of ``isotropic'' elements:``Circles placed arbitrarily on a plane'', he says, ``are therefore not totally independent among themselves, as could be thought at first sight, they have ideally two imaginary points in common at infinity'' ({[}252{]}, v. I, p.~48). Later, he introduces in the same way the ``umbilical'', an imaginary conic at infinity common to all spheres ({[}252{]}, v. I, p.~370); and if he does not speak particularly of the isotropic generators of the sphere, at least he underlines explicitly the existence of rectilinear generators, real or imaginary, for all quadrics (ibid., p.~371); \(^{15}\) notions that his followers (notably

Plücker and Chasles), even more than he, make great use of, in particular in the study of ``focal'' properties of conics and quadrics.

\begin{enumerate}
\def\labelenumi{\Alph{enumi})}
\setcounter{enumi}{2}
\item
  The notions of point transformation and of composition of transformations are, they also, formulated in a general way and introduced systematically as means of proof. Apart from displacements and projections, only certain particular transformations were known up to that point: certain plane projective transformations of the type \(x' = a / x\) , \(y' = y / x\) , used by La Hire and Newton, the ``affinity'' \(x' = ax\) , \(y' = by\) of Clairaut and Euler ({[}108 a{]}, (1), v. IX, chap.~XVIII), and finally a few particular quadratic transformations, with Newton again, Maclaurin and Braikenridge. Monge, in his Géométrie descriptive, shows all the use that can be made of plane projections in 3 dimensional geometry. With Poncelet, one of the systematic procedures for proofs, used to surfeit, consists in reducing by projection properties of the conic to those of the circle (a method already used on occasion by Descartes and Pascal); and in order to go from a quadric to a sphere, he invents the first example of a projective transformation of space, the ``homology'' ({[}252{]}, v. I, p.~357); finally it is he also who introduces the first examples of birational transformations of a curve into itself. In 1827, Möbius ({[}223{]}, v. I, p.~217) (and independently Chasles in 1830 ({[}60 b{]}, p.~695)), define the most general linear projective transformations; at the same time appear inversion and other types of quadratic transformations, whose study will inaugurate the theory of birational transformations, which will develop in the second half of the XIXth century.
\item
  The notion of duality appears fully in the open and is found consciously linked to the theory of bilinear forms. The theory of poles and polars with respect to conics, which, since Apollonius, had only made some progress with Desargues and La Hire, is extended to quadrics by Monge, who, as well as his pupils, sees the possibility of transforming by this means known theorems into new results. \(^{16}\) But it is again due to Poncelet that the credit returns for having built up these remarks into a general method in the theory of transformations ``by reciprocal polars'', and of having made it into a particularly efficacious tool of discovery. A little later, notably with Gergonne, Plücker, Möbius and Chasles, the general notion of duality frees itself from the link with quadratic forms, still too close with Poncelet. In particular, Möbius, while examining the various possibilities for duality in space of 3 dimensions (defined by a bilinear form), discovers in 1833 duality with respect to an alternating bilinear form ({[}223{]}, v. I, pp.~489-515), \(^{17}\) especially studied, in the XIXth century, under the form of the theory of ``linear complexes'' and developed in relation to the ``geometry of straight lines'' and the ``Plücker co-ordinates'' introduced by Cayley, Grassmann and Plücker around 1860.
\item
  From the beginnings of projective geometry, the intensive study of the properties of classical geometry in their relationships with projective space had rapidly led to dividing them into ``projective properties'' and ``metric properties''; and it is no doubt not exaggerated to see in this separation one of the clearest manifestations, at this time, of what was to become the modern notion of structure. But Poncelet, who is the first to introduce this distinction and this terminology, is already conscious of what links these two types of properties; and, tackling in the Traité the problems concerning angles, whose properties ``seem not to be part of those that we have called projective \ldots, they follow on nonetheless in such a simple way'', he says, ``from the principles that form the basis {[}of this work{]} \ldots, that I do not believe that any other geometrical theory could lead to it in a manner which is both more direct and simpler. It will not be astonishing, if one considers that projective properties of figures are necessarily the most general among those that can belong to them; in such a way that they must encompass, as simple corollaries, all the other properties or particular relations of the area'' ({[}252{]}, v. I, p.~248). Truth to tell, after this declaration, it is a little surprising to see him tackle questions about angles in a very roundabout way, by linking them to focal properties of conics, instead of bringing in directly the cyclic points; and in fact, it is only 30 years later that Laguerre (still a pupil at the École Polytechnique) gave the expression for a right angle by means of the cross-ratio of these straight lines and isotropic straight lines with the same origin ({[}192{]}, v. II, p.~13). Finally, with Cayley ({[}58{]}, v. II, pp.~561-592) the fundamental idea that ``metric'' properties of a plane figure are none other than the ``projective'' properties of the figure extended by the cyclic points becomes clearly expressed --- a decisive way station towards the ``Erlangen programme''.
\item
  Non-Euclidean hyperbolic geometry, which will come to light around 1830, remains at first a little apart from the movement of which we are tracing the main features. Arising out of preoccupations of an essentially logical order touching on the foundations of classical geometry, this new geometry is presented by its inventors \(^{18}\) in the same axiomatic and ``synthetic'' form as Euclid's geometry, and without any link to projective geometry (whose introduction following the classical model seems even to be excluded a priori, since the notion of unique parallel disappears in this geometry); it is no doubt because of that that it hardly attracts, for a long time, any interest from the French, German or English schools of projective geometry. Also, when Cayley, in the fundamental memoir quoted earlier ({[}58{]}, v. II, pp.~561-
\end{enumerate}

\begin{enumerate}
\def\labelenumi{\arabic{enumi})}
\setcounter{enumi}{591}
\tightlist
\item
  has the idea of replacing cyclic points (considered as a ``tangentially degenerate'' conic) by an arbitrary conic (that he calls ``absolute''), he is not thinking at all of linking this idea with the geometry of Lobatschevsky-Bolyai, even though he indicates how his conception leads to new expressions for the ``distance'' between two points, and that he mentions its links with spherical geometry. The situation changes around 1870, when the non-Euclidean geometries, following the diffusion of the works of Lobatschevsky, and the publication of the works of Gauss and the inaugural lecture of Riemann, come to the fore of mathematical news. Following the path traced by Riemann, Beltrami, without knowing of the work of Cayley, finds again in 1868 expressions for distance as given by this latter, but in quite another context, by considering the interior of a circle as an image of a surface with constant curvature, in which the geodesics are represented by straight lines {[}18 a{]}; it is Klein, who two years later, makes (independently of Beltrami) a synthesis of these various points of view, which he completes by the discovery of elliptic non-Euclidean space ({[}182{]}, v. I, pp.~254-305). \(^{19}\)
\end{enumerate}

\begin{enumerate}
\def\labelenumi{\Alph{enumi})}
\setcounter{enumi}{6}
\tightlist
\item
  In the second half of the period which we are considering here, a period of critical reflection is inserted, during which the supporters of ``synthetic'' geometry, not content with having banished co-ordinates from their proofs, pretend to do without real numbers even in their axioms for geometry. The principal representative of this school is von Staudt, who succeeds essentially in bringing about this tour de force {[}325{]}, very admired in his time and even well into the XXth century; and if today the same importance is not attributed to ideas of this order, whose possibilities for fruitful application are revealed to be fairly thin, it must nonetheless be recognised that the efforts of von Staudt and his disciples contributed to clarifying the role of real or complex ``scalars'' in classical geometry, and to introducing in that very way the modern conception of geometries over an arbitrary base field.
\end{enumerate}

Around 1860, ``synthetic'' geometry was at its peak, but the end of its reign was approaching at great speed. Having remained heavy and ungracious during the whole of the XVIIIth century, analytic geometry, in the hands of Lamé, Bobillier, Cauchy, Plücker and Möbius, acquires at last the elegance and conciseness that will allow it to fight on equal terms with its rival. Especially, after about 1850, the ideas of group and invariant, formulated at last in a precise way, invade little by little the scene, and it is seen that the theorems of classical geometry are none other than the expression of identical relations between invariants or covariants of the group of similarities, \(^{20}\) in the same way that those of projective geometry express the identities (or ``syzygies'') between covariants of the projective group. It is this thesis that is magisterially stated by F. Klein in the famous ``Erlangen programme'' ({[}182{]}, v. I, pp.~460-497), where he advocates the abandonment of sterile controversies between the ``synthetic'' tendency and the ``analytical'' tendency; if, he says, the accusation levelled against this latter of giving a privileged role to an arbitrary system of axes ``was most often justified in what concerns the defective way in which the method of co-ordinates was used previously, it founders when it is a case of a rational application of this method . . . The domain of spatial intuition is not forbidden by the analytic method . . .'', and he underlines that ``one must not underestimate the advantage that a well adapted formalism brings to subsequent research, in that it goes in advance so to say of thought'' (loc. cit., pp.~488-490).

One concludes thus with a rational and ``structural'' classification of the theorems of ``geometry'' according to the group from which they arise: the linear group for projective geometry, the orthogonal group for metric questions, the symplectic group for ``linear complex'' geometry. But under this merciless clarity, classical geometry --- with the exception of algebraic geometry and differential geometry, \(^{21}\) henceforth constituted as autonomous sciences --- flares suddenly and loses all its brightness. Already the generalisation of methods based on the use of transformations had made the development of new theorems somewhat mechanical: ``Today,'', says Chasles in 1837 in his Aperçu historique, ``every one can put himself forward, take an arbitrary known truth, and submit it to various general transformation principles; he will draw out other truths, different or more general; and these will be susceptible to similar operations; so that one will be able to multiply, almost to infinity, the number of new truths deduced from the first \ldots{} So anyone can who wishes, in the present state of the science, generalise and create in Geometry; genius is no longer indispensable in order to add a stone to the structure'' ({[}60 b{]}, pp.~268-269). But the situation becomes much clearer with progress in the theory of invariants, which succeeds at last (at least for the ``classical'' groups) in formulating general methods allowing in principle the writing down of all the algebraic covariants and all their ``syzygies'' in a purely automatic way; a victory that, at the same time, signals the death, as a field of research, of the classical theory of invariants itself, and of ``elementary'' geometry, \(^{22}\) which has become in practice a simple dictionary for it. No doubt, nothing allows the foresight a priori, among the infinity of ``theorems'' that can thus be rolled out at will, as to which will be those whose statement, in the appropriate geometric language, will have a simplicity and an elegance comparable to those of the classical results, and there remains there a restricted domain where numerous amateurs continue to work with pleasure (the geometry of the triangle, the tetrahedron, of algebraic curves and surfaces of low degree, etc.). But for the professional mathematician, the mine is worked out, since there are no longer there any structural problems, susceptible of reverberating on other parts of mathematics; and this chapter of the theory of groups and invariants can be considered as closed until further notice. \(^{23}\)

Thus, after the Erlangen programme, Euclidean and noneuclidean geometries, from a purely algebraic point of view, became simple languages, more or less convenient, for expressing the results of the theory of bilinear forms, whose progress went in tandem with that of the theory of invariants. \(^{24}\) All that concerns the notion of rank of a bilinear form and the relationships between these forms and linear transformations is definitively clarified by the works of Frobenius ({[}119{]}, v. I, pp.~343-405). It is also due to Frobenius that we have the canonical expression of an alternating form on a free Z-module ({[}119{]}, v. I, pp.~482-544); however, the left symmetric determinant had already appeared with Pfaff, at the beginning of the century, in connection with the reduction of differential forms to a normal form; Jacobi, who, in 1827, takes up this problem again ({[}171{]}, v. IV, pp.~17-29), knows that a left symmetric determinant of odd order is null, and it is he who constructs the expression for the pfaffian and shows that it is a factor of the left symmetric determinant of even order; but he had not realised that this latter is the square of the pfaffian, and this point was only established by Cayley in 1849 ({[}58{]}, v. I, pp.~410-413). The notion of the symmetric bilinear form associated with a quadratic form is the most elementary case of the process of ``polarisation'', one of the fundamental tools of the theory of invariants. Under the name of ``scalar product'', this notion will know an immense prosperity, first with the vulgarisers of the ``vector calculus'', then, starting in the XXth century, thanks to the unsuspected generalisation that is brought to it by the theory of Hilbert spaces (see p.~212). It is also this latter theory that will bring to light the notion of the adjoint of an operator (which beforehand had hardly come out except in the theory of linear differential equations, and, in tensor calculus, by the waltz of the co- and contravariant indices under the baton of the metric tensor); it is that which finally will give all its character to the notion of hermitian form introduced first by Hermite in 1853 in connection with his arithmetic research ({[}159{]}, v. I, p.~237), but having remained on the edge of the great mathematical streams until around 1925 and the applications of complex Hilbert spaces to quantum theories.

The study of the orthogonal group and the group of similarities --- clearly conceived and treated as such since the middle of the XIXth century, and having become the heart of the theory of quadratic forms --- as well as that of the other ``classical'' groups (linear group, symplectic group and unitary group), takes up on the other hand a greater and greater importance. We can only mention here the essential role played by these groups, in the theory of Lie groups and differential geometry on the one hand, the arithmetic theory of quadratic forms (see for example {[}285{]} and {[}100{]}) on the other; \(^{25}\) it is to this circumstance, as well as to the extension of the concept of duality to the most diverse questions that is due the fact that there is hardly any modern mathematical theory where bilinear forms do not crop up in some form or other. We must in any case note that it is the study of rotation groups (in three dimensions) that led Hamilton to the discovery of quaternions {[}145 a{]}; this discovery is generalised by W. Clifford who, in 1876, introduces the algebras that bear his name, and proves that they are tensor products of quaternion algebras, or quaternion algebras and a quadratic extension ({[}65{]}, pp.~266-276). Found again four years later by Lipschitz {[}205 b{]}, who uses them to give a parametric representation of orthogonal transformations in n variables (generalising those that Cayley had obtained for n = 3 ({[}58{]}, v. I, pp.~123-126) and n = 4 (v. II, pp.~202-215) by the theory of quaternions), these algebras, and the notion of ``spinor'' that derives from them ({[}52 b{]} and {[}62 a{]}), would also know a great vogue in modern times by virtue of their use in quantum theories.

It remains for us finally to say a word about the evolution of ideas which led to the almost total abandonment of all restrictions on the ring of scalars in the theory of sesquilinear forms --- a common tendency in all modern algebra, but which has perhaps manifested itself here earlier than elsewhere. We have already pointed out the fruitful introduction of geometry over the field of complex numbers (which, besides, during the whole of the XIXth century, did not go without a perpetual and sometimes dangerous confusion between this geometry and real geometry); the clarity here comes especially from the axiomatic studies of the end of the XIXth century on the foundations of geometry {[}163 c{]}. In the course of this research, Hilbert and his emulators, notably, on examining the relationships between the various axioms, were led to construct appropriate counterexamples, where the ``base field'' (commutative or not) possessed more or less pathological properties, and they thus accustomed mathematicians to ``geometries'' of a quite new type. From an analytical point of view, Galois had already considered linear transformations where coefficients and variables took values in a finite prime field ({[}123{]}, p.~145); in developing these ideas, Jordan {[}174 a{]} is led in a natural way to consider the classical groups over these fields, groups whose intervention manifests itself in varied mathematical domains. Dickson, around 1900, extends the research of Jordan to all finite fields, and more recently, it has become apparent that a large part of the theory of Jordan-Dickson can be extended to an absolutely arbitrary ``base field''; this is due essentially to general properties of isotropic vectors and to Witt's theorem, which, trivial in the classical cases, were only established for an arbitrary base field in 1936 {[}337 a{]}. \(^{26}\)

But in thus pushing the study of sesquilinear forms towards an ``abstraction'' which is always greater, it has turned out to be extremely suggestive to keep as is the terminology which, in the case of space of 2 or 3 dimensions, came out of classical geometry, and to extend it to the n-dimensional case and even to infinite dimensional spaces. Passed over in its role as an autonomous and living science, classical geometry is thus transfigured into a universal language of contemporary mathematics, with a suppleness and a usefulness that are incomparable.

